\documentclass[openany, twoside]{book}

% Language setting
%% Replace `english' with e.g. `spanish' to change the document language
\usepackage[english]{babel}

% Set page size and margins
%% Replace `letterpaper' with `a4paper' for UK/EU standard size
\usepackage[letterpaper,top=2cm,bottom=2cm,left=3cm,right=3cm,marginparwidth=1.75cm]{geometry}

% Properly format custom page numbering system
%% Roman numerals for preface pages and Arabic numbers for every other page
\usepackage{emptypage}

\usepackage{fancyhdr}
\pagestyle{fancy}

\renewcommand{\chaptermark}[1]{\markboth{#1}{}}

\fancypagestyle{frontmatter}{
    \fancyhf{}
    \fancyfoot[EL]{\textbf{\thepage}}
    \fancyfoot[OR]{\textbf{\thepage}}
    \renewcommand{\headrulewidth}{0pt}
}

\fancypagestyle{mainmatter}{
    \fancyhf{}
    \fancyfoot[EL]{\textbf{\thepage}\ /\ Chapter \thechapter}
    \fancyfoot[OR]{\leftmark\ /\ \textbf{\thepage}}
    \renewcommand{\headrulewidth}{0pt}
}

\makeatletter
\fancypagestyle{plain}{
    \fancyhf{}
    \if@mainmatter
        \fancyfoot[EL]{\textbf{\thepage}\ /\ Chapter \thechapter}
        \fancyfoot[OR]{\leftmark\ /\ \textbf{\thepage}}
    \else
        \fancyfoot[EL]{\textbf{\thepage}}
        \fancyfoot[OR]{\textbf{\thepage}}
    \fi
    \renewcommand{\headrulewidth}{0pt}
}
\makeatother

% Other packages
\usepackage{amsmath}
%\usepackage{amsthm}
\usepackage{amssymb}

\usepackage{ntheorem}
\theorembodyfont{\normalfont}

\usepackage{graphicx}
\usepackage[colorlinks=true, allcolors=blue]{hyperref}

\usepackage{multicol}
\usepackage{multirow}
\usepackage{array}
\usepackage{tabularx}
\usepackage{tabularray}

\usepackage{verbatim}

\usepackage{indentfirst}
\usepackage[autostyle=true]{csquotes}
\MakeOuterQuote{"}

\usepackage{titlesec}
\titleformat{\chapter}[block]
  {\normalfont\Huge\bfseries}
  {\thechapter \hspace{1.5mm}}
  {0pt}
  {\LARGE\bfseries /\ }

\titlespacing*{\chapter}{0pt}{10pt}{10pt}

% Proper figure/table numbering, wrapping, and captions
\usepackage{chngcntr}
\counterwithin{figure}{chapter}
\counterwithin{table}{chapter}
\usepackage{wrapfig}

\usepackage{caption}
\captionsetup{labelfont=bf, labelsep=space}

% Diagrams and geometric figures
\usepackage{tikz}
\usetikzlibrary{intersections}

% Theorem-Like Environments
\newcounter{theoremcounter}
\counterwithin{theoremcounter}{chapter}
\NewDocumentEnvironment{thm}{O{} b}{%
    \refstepcounter{theoremcounter}%
    \begin{minipage}{\textwidth}
    \par\vspace{1em} \noindent \textbf{Theorem \thetheoremcounter} \textit{#1}%
    \hspace*{0.5mm} \rule{0.25\linewidth}{0.4pt}%
    \\
    \indent #2%
    \\[-5pt]
    \rule{0.5\linewidth}{0.4pt}%
    \end{minipage}
    }{%
        \medskip%
    }

\NewDocumentEnvironment{corollary}{O{} b}{%
    \refstepcounter{theoremcounter}%
    \begin{minipage}{\textwidth}
    \par\vspace{1em} \noindent \textbf{Corollary \thetheoremcounter} \textit{#1}%
    \hspace*{0.5mm} \rule{0.25\linewidth}{0.4pt}%
    \\
    \indent #2%
    \\[-5pt]
    \rule{0.5\linewidth}{0.4pt}%
    \end{minipage}
    }{%
        \medskip%
    }

\newcounter{postulatecounter}
\NewDocumentEnvironment{post}{O{} b}{%
  \refstepcounter{postulatecounter}%
  \begin{minipage}{\textwidth}
  \par\vspace{1em} \noindent \textbf{Postulate \thepostulatecounter} \textit{#1}%
    \hspace*{0.5mm} \rule{0.25\linewidth}{0.4pt}%
    \\
    \indent #2%
    \\[-5pt]
    \rule{0.5\linewidth}{0.4pt}%
    \end{minipage}
}{%
  \medskip%
}

\begin{comment}
Format for theorems:
\begin{thm}[AAS Theorem]
    If two triangles have two corresponding sides congruent and a non-included angle congruent, they are congruent.
\end{thm}

Format for postulates:
\begin{post}[xyz Postulate]
    If two triangles have two corresponding sides congruent and a non-included angle congruent, they are congruent.
\end{post}
\end{comment}

\newcounter{proofstep}
\newcommand{\step}[2]{%
  \stepcounter{proofstep}%
  \theproofstep. & #1 & #2 \\%
}

\begin{comment}
Format for proofs:

\setcounter{proofstep}{0} 
\medskip \noindent
\begin{tabular}{r|p{4cm}|p{8cm}}
  \textbf{No.} & \textbf{Statements} & \textbf{Reasons} \\ \hline
  \step{blahblah}{blahblahblahblah}
  \step{blahblah}{blahblahblahblah}
  \step{blahblah}{blahblahblahblah}
  \step{blahblah}{blahblahblahblah}
\end{tabular}
\smallskip
\hfill Q.E.D. \hspace{0.5cm}
\medskip
\end{comment}

\newcounter{constructioncounter}
\NewDocumentEnvironment{construct}{O{} b}{%
    \refstepcounter{constructioncounter}%
    \begin{minipage}{\textwidth}
    \par\vspace{1em} \noindent \textbf{Construction \theconstructioncounter} \textit{#1}%
    \hspace*{0.5mm} \rule{0.25\linewidth}{0.4pt}%
    \\
    \indent #2%
    \\[-5pt]
    \rule{0.5\linewidth}{0.4pt}%
    \end{minipage}
    }{%
        \medskip%
    }

 % Other stuff in preamble
\raggedbottom

% DOCUMENT ——————————————————————————————
\begin{document}

\begin{titlepage}
    \vspace*{1mm}
    {\Huge \bfseries Geometry}
    \vspace*{\fill}
    \begin{flushright}
        {\Large \bfseries Textbook by Rafan Rayan} \\
        \smallskip
        Prerelease Version 5
    \end{flushright}
    \vspace*{1mm}
\end{titlepage}

\thispagestyle{empty}
\vspace*{\fill}

\begin{center}
    \small This document was typeset using \LaTeX.
\end{center}

\clearpage

\frontmatter
\pagestyle{frontmatter}

\section*{Preface}

Geometry is a foundational branch of mathematics focused on the properties, measurements, and relationships of points, lines, angles, surfaces, and solids. It provides the logical framework for understanding the size, shape, and relative position of figures in both two-dimensional planes and three-dimensional spaces, allowing us to quantify the physical world.

This field is crucial because it underpins countless practical applications in daily life, science, and technology. For instance, geometry is essential for architects and engineers to design safe structures, maps the principles behind computer graphics and video game development, and allows medical professionals to utilize advanced imaging technologies like CT scans.

Beyond applied sciences, studying geometry trains the brain in spatial reasoning and deductive logic. It teaches you how to analyze problems, construct logical proofs, and visualize abstract concepts. Whether you are calculating the area of a room, navigating using a GPS, or launching a satellite, geometry provides the spatial rules that make it all possible.

\section*{Contents}

\subsection*{1 \indent The Fundaments of Geometry}
\noindent \nameref{sec:1/Logic} \dotfill \pageref*{sec:1/Logic} \\
\noindent \nameref{sec:1/Euclid and Euclidean Geometry} \dotfill \pageref*{sec:1/Euclid and Euclidean Geometry} \\
\noindent \nameref{sec:1/Fundamental Figures} \dotfill \pageref*{sec:1/Fundamental Figures} \\
\noindent \nameref{sec:1/Proofs} \dotfill \pageref*{sec:1/Proofs} \\
\noindent \nameref{sec:1/Useful Properties, Postulates, Definitions, and Theorems} \dotfill \pageref*{sec:1/Useful Properties, Postulates, Definitions, and Theorems}

\subsection*{2 \indent Parallelism}
\noindent \nameref{sec:1/Logic} \dotfill \pageref*{sec:1/Logic} \\

\subsection*{3 \indent Congruent Triangles}
\noindent \nameref{sec:1/Logic} \dotfill \pageref*{sec:1/Logic} \\

\subsection*{4 \indent Inequalities}
\noindent \nameref{sec:1/Logic} \dotfill \pageref*{sec:1/Logic} \\

\subsection*{5 \indent Quadrilaterals}
\noindent \nameref{sec:1/Logic} \dotfill \pageref*{sec:1/Logic} \\

\subsection*{6 \indent Similarity}
\noindent \nameref{sec:1/Logic} \dotfill \pageref*{sec:1/Logic} \\

\subsection*{7 \indent Right Triangles}
\noindent \nameref{sec:1/Logic} \dotfill \pageref*{sec:1/Logic} \\

\subsection*{8 \indent Circles}
\noindent \nameref{sec:8/Defining a Circle} \dotfill \pageref*{sec:8/Defining a Circle} \\

\subsection*{9 \indent Solids}
\noindent \nameref{sec:1/Logic} \dotfill \pageref*{sec:1/Logic} \\

\subsection*{10 \indent Coordinate Geometry}
\noindent \nameref{sec:1/Logic} \dotfill \pageref*{sec:1/Logic} \\

\subsection*{11 \indent Transformations}
\noindent \nameref{sec:1/Logic} \dotfill \pageref*{sec:1/Logic} \\

\subsection*{12 \indent Constructions}
\noindent \nameref{sec:12/Tools in Constructions} \dotfill \pageref*{sec:12/Tools in Constructions} \\

\mainmatter
\pagestyle{mainmatter}

\include{01_the-fundaments-of-geometry}
\include{02_parallelism}
\include{03_congruent-triangles}
\include{04_inequalities}
\include{05_quadrilaterals}
\include{06_similarity}
\include{07_right-triangles}
\include{08_circles}
\include{09_solids}
\include{10_coordinate-geometry}
\include{11_transformations}
\include{12_constructions}

\end{document}
